\documentclass[12pt]{article}
\usepackage[utf8]{inputenc}
\usepackage[T1]{fontenc}
\usepackage[margin=1in]{geometry}
\usepackage{setspace}
\usepackage{xcolor}
\usepackage[colorlinks=true, urlcolor=blue!60!black, linkcolor=blue!60!black]{hyperref}
\setlength{\parindent}{2em}
\setlength{\parskip}{1.25em}
\renewcommand{\baselinestretch}{1.0}
\begin{document}
\begin{center}
{\LARGE \textit{New York Times Co. v. Sullivan}, 376 U.S. 254 (1964)\par}
\vspace{0.8em}
{\large Case Brief --- Four-Page Edition\par}
\end{center}
\vspace{1em}

\section{Memory Jogger}

An Alabama jury awarded a Montgomery city commissioner \$500,000 against the \textit{New York Times} over a civil-rights advertisement with minor errors; the Court unanimously reversed, holding that a public official may recover for a defamatory falsehood about his official conduct only by proving ``actual malice''---knowing or reckless falsity (\href{https://supreme.justia.com/cases/federal/us/376/254/}{Justia}).

\section{Detailed Case Facts}

On March 29, 1960, the \textit{Times} published ``Heed Their Rising Voices,'' a full-page advertisement paid for by the Committee to Defend Martin Luther King and the Struggle for Freedom in the South. It described a ``wave of terror'' against Southern demonstrators and appealed for funds; four Alabama clergymen listed as endorsers testified that they had never authorized the use of their names. Only about 394 copies circulated in Alabama. L.B. Sullivan, the Montgomery commissioner who supervised the police, was never named, but he sued over passages with conceded minor errors: students sang the National Anthem, not ``My Country, 'Tis of Thee''; the dining hall was never padlocked; police did not ``ring'' the campus; and Dr. King had been arrested four times, not seven. The \textit{Times} answered his retraction demand by saying it was ``puzzled'' how the text reflected on him.

\section{Procedural History}

Sullivan offered no proof of pecuniary loss. The trial court instructed that the statements were ``libelous per se,'' so falsity, malice, and general damages were presumed, and the jury awarded the full \$500,000 against all defendants. The Supreme Court of Alabama affirmed, holding that the First Amendment does not protect libel, that the Fourteenth Amendment does not reach private suits, and that criticism of the police necessarily concerned their supervisor.

\section{Judicial Votes}

Reversed and remanded, \textbf{9--0} (decided March 9, 1964).
\begin{itemize}
\item \textbf{Majority:} Justice Brennan, joined by Chief Justice Warren and Justices Clark, Harlan, Stewart, and White.
\item \textbf{Concurring:} Justice Black, joined by Justice Douglas; Justice Goldberg (in the result), joined by Justice Douglas.
\end{itemize}

\section{Holding}

The First and Fourteenth Amendments bar a public official from recovering damages for a defamatory falsehood relating to his official conduct unless he proves with ``convincing clarity'' that the statement was made with actual malice---knowledge of falsity or reckless disregard of whether it was false. On independent review, the evidence was constitutionally insufficient to show actual malice or that the impersonal criticism of government was ``of and concerning'' Sullivan.

\section{Analysis of Opinions}

\textbf{Majority (Brennan).} \emph{First}, enforcing common-law tort rules in state court is state action, and a paid editorial advertisement does not forfeit protection. \emph{Second}, libel enjoys no ``talismanic immunity.'' Erroneous statements are inevitable in free debate and need ``breathing space.'' The Sedition Act of 1798 controversy revealed the ``central meaning of the First Amendment''; civil damages here were 1,000 times Alabama's maximum criminal fine, with no protection against cascading suits. A truth defense fails because critics would ``steer far wider of the unlawful zone.'' By analogy to officials' immunity in \textit{Barr v. Matteo}, the citizen-critic deserves a privilege defeasible only by actual malice. \emph{Third}, the Court independently examined the record: there was no proof against the clergymen, and the \textit{Times}'s failure to check its own files was at most negligence. Treating criticism of government as libel of officials revived seditious libel through ``legal alchemy.''

\textbf{Black (concurring).} The Constitution confers an absolute immunity to criticize officials; with eleven Alabama suits seeking millions against the \textit{Times} and CBS, actual malice offers too little protection.

\textbf{Goldberg (concurring in the result).} He too would grant an absolute privilege for criticism of official conduct, since juries easily conflate falsity with malice, while leaving purely private defamation actionable.

\textbf{Later.} \textit{Curtis Publishing v. Butts} extended the rule to public figures; \textit{Gertz} let states use lower fault standards for private plaintiffs; \textit{Hustler v. Falwell} applied it to emotional distress; \textit{Bose} reinforced independent review.

\section{Examples: Future Applications}

\begin{enumerate}
\item \textbf{Same-Side (No Liability): Careless Op-Ed.} A paper runs an op-ed with wrong dates about an Attorney General without checking the docket---at most negligence.
\item \textbf{Same-Side (No Liability): The Sheriff and the Jail.} A false charge that unnamed jail guards beat inmates is not ``of and concerning'' the elected sheriff.
\item \textbf{Opposite-Side (Malice Shown): Purposeful Avoidance.} A blog relies on a tipster known to fabricate and refuses to check public audit logs refuting the claim against a police chief (\textit{Harte-Hanks}).
\item \textbf{Opposite-Side (Malice Not Required): Private Mechanic.} A private mechanic falsely accused of overbilling may recover actual damages on negligence under \textit{Gertz}.
\item \textbf{Fence-Sitter: One-Time Speaker.} A business owner who testified once at a zoning hearing is falsely accused of bribery. Whether she is a limited-purpose public figure is uncertain.
\end{enumerate}

\section{Critique}

\textbf{Competing critiques.} Originalists, echoing Justice Thomas, argue that neither the First nor the Fourteenth Amendment displaced state libel law, noting that Madison contemplated state remedies for injured reputations. Absolutists in the tradition of Black and Goldberg answer that intrusive discovery into editors' states of mind chills underfunded speakers.

\textbf{The digital ecosystem.} Pragmatists add that where falsehoods cost little to spread, actual malice may under-deter deliberate lies while leaving wrongly maligned limited-purpose public figures without redress.

\section{Key Quotations}

\begin{quote}
``The constitutional guarantees require, we think, a federal rule that prohibits a public official from recovering damages for a defamatory falsehood relating to his official conduct unless he proves that the statement was made with `actual malice'---that is, with knowledge that it was false or with reckless disregard of whether it was false or not.''
\end{quote}
\noindent --- Justice Brennan, for the Court, 376 U.S. at 279--80.

\begin{quote}
``Thus, we consider this case against the background of a profound national commitment to the principle that debate on public issues should be uninhibited, robust, and wide-open, and that it may well include vehement, caustic, and sometimes unpleasantly sharp attacks on government and public officials.''
\end{quote}
\noindent --- Justice Brennan, for the Court, 376 U.S. at 270.

\begin{quote}
``Whether or not a newspaper can survive a succession of such judgments, the pall of fear and timidity imposed upon those who would give voice to public criticism is an atmosphere in which the First Amendment freedoms cannot survive.''
\end{quote}
\noindent --- Justice Brennan, for the Court, 376 U.S. at 278.

\end{document}
